\section{Sharp memory on complex projective orbits}\label{sec:projective}
Write $r_q=\sqrt{1-1/q}$ and associate to a rank-one orthogonal
projector $P$ the unit vector
\begin{equation}\label{eq:pure-orbit}
 Y(P)=\frac{P-I/q}{r_q}\in\Herm(q),\qquad
 \mathcal K_q=\conv\{Y(P):\rank P=1\}.
\end{equation}
The body $\mathcal K_q$ consists of normalized centered density matrices.
It is invariant under conjugation and contains zero. Its extreme orbit
is $\CP^{q-1}$. Fix an orthonormal real basis $F_1,\ldots,F_D$ of
$\Herm(q)$ and use its coordinate queries in \eqref{eq:law}.
One may take as seeds the $q^2$ rank-one projectors onto
$e_i$, $(e_i+e_j)/\sqrt2$ and $(e_i+i e_j)/\sqrt2$, $i<j$.
Their centered matrices affinely span $\Herm(q)$. All seed and query
alphabets are fixed when $q$ is fixed.

\subsection{A projective enclosure with quadratic loss}
For rank-one projectors use the distance
$d(P,Q)=\sqrt{1-\tr(PQ)}$, the sine of their principal angle.
\begin{lemma}[Finite projective synthesis]\label{lem:projective-net}
For $0<h\le1$, there is an $h$-net $\mathcal V_h$ of rank-one
projectors with at most
\begin{equation}\label{eq:net-cardinality}
 q\bigl(2\lceil2\sqrt{2(q-1)}/h\rceil+1\bigr)^{2(q-1)}
\end{equation}
elements, and the entries of all these projectors can be chosen in
$\mathbb Q+i\mathbb Q$. If $qh^2<1$, then
\begin{equation}\label{eq:projective-enclosure}
 (1-qh^2)\mathcal K_q\subseteq
 \conv\{Y(P):P\in\mathcal V_h\}\subseteq\mathcal K_q.
\end{equation}
\end{lemma}
\begin{proof}
Every complex line has a representative with a largest coordinate
scaled to one and the real and imaginary parts of all other coordinates
in $[-1,1]$. In each of the $q$ charts round these $2(q-1)$ real
coordinates to the grid of step $1/M$, where
$M=\lceil2\sqrt{2(q-1)}/h\rceil$. If $v,w$ are the original and
rounded representatives, their norms are at least one and
$\norm{v-w}\le\sqrt{2(q-1)}/(2M)$. Normalization changes their
distance by at most a factor two, and the sine distance is at most
the distance of these normalized representatives. Thus the stated
set is an $h$-net. The projector $ww^*/\norm w^2$ has rational real
and imaginary entries, and the grid count is \eqref{eq:net-cardinality}.

For a nonzero traceless Hermitian $H$, let $\lambda_1$ and $\lambda_q$
be its largest and smallest eigenvalues. A net projector within $h$
of a top eigenline satisfies
\[
 \tr(HP)\ge\lambda_1-(\lambda_1-\lambda_q)h^2
                \ge(1-qh^2)\lambda_1.
\]
Indeed $-\lambda_q\le(q-1)\lambda_1$ by trace zero. The support
function of $\mathcal K_q$ at $H$ is $\lambda_1/r_q$.
The support-function criterion for convex containment proves
\eqref{eq:projective-enclosure}. It applies to every direction, not
just a finite set of tested queries.
\end{proof}

\begin{theorem}[A finite-alphabet exact upper machine]
\label{thm:projective-upper}
For every finite unitary command alphabet and $0<\rho\le1/10$, the
experiment on \eqref{eq:pure-orbit} has an exact machine of width
$O_q(N^{q-1})$. For each $N$, all command transition matrices can be
chosen independent of the epoch. Each row uses at most $q^2$ successors.
The final decoder may depend on $N$.
\end{theorem}
\begin{proof}
Take $h=(8qN)^{-1/2}$ and $\eta=1-1/(8N)$ in
Lemma~\ref{lem:projective-net}. Let $V$ be the resulting finite set of
centered vectors, and put $a_t=\eta^{-t}/4$. A label $i$ represents
$a_tY(P_i)$ at epoch $t$. This is a proof description of a label;
no matrix is stored in the register.

Since $\rho/a_0=4\rho\le2/5<\eta$ and $0\in\mathcal K_q$,
any seed $\rho Y(P)$ is a convex combination of $a_0V$.
For each label $i$ and each command $U$, the vector
$\eta UY(P_i)U^*$ lies in $\conv V$. Choose a convex representation
and use its weights as the stochastic transition row. The choices depend
on the command and label but not on $t$. Multiplying by $a_{t+1}$
shows that the expected represented vector is updated by exact
conjugation, since $a_{t+1}\eta=a_t$.

Induction gives expected vector $\rho Y(U_wPU_w^*)$. The terminal
column for query $j$ is $a_N\tr(F_jY(P_i))$. It is legal because
$|\tr(F_jY(P_i))|\le1$ and Bernoulli's inequality gives
$\eta^N\ge7/8$, so $a_N\le2/7$. Sampling a binary answer with this
mean proves every conditional probability in \eqref{eq:law} exactly.
Carath\'eodory's theorem in real dimension $D=q^2-1$ permits $q^2$
successors. The number in \eqref{eq:net-cardinality} is
$O_q(h^{-2(q-1)})=O_q(N^{q-1})$.
\end{proof}
The upper argument uses a quadratic support loss along the projective
orbit. Applying a net on the full ambient sphere would instead give
the larger exponent $(q^2-2)/2$.

\subsection{Fixed finite algebraic alphabets}
For $1\le j<q$ let $R_j,D_j\in\SU(q)$ act as the identity outside
the adjacent coordinate pair $(j,j+1)$ and on that pair as
\begin{equation}\label{eq:adjacent-gates}
 R=\frac15\begin{pmatrix}3&-4\\4&3\end{pmatrix},\qquad
 D=\begin{pmatrix}\zeta&0\\0&\overline\zeta\end{pmatrix},\qquad
 \zeta=\frac{3+4i}{5}.
\end{equation}
The alphabet has $4(q-1)+1$ symbols:
$\{I,R_j^{\pm1},D_j^{\pm1}:1\le j<q\}$.
Assign probability $1/2$ to $I$ and $1/[8(q-1)]$ to each other symbol.

\begin{proposition}[Qualitative group expansion]\label{prop:algebraic-gates}
This command law has a norm gap on $L^2_0(\SU(q))$ for every fixed $q$.
Its entries and probabilities are algebraic. No numerical gap value is
asserted for this family.
\end{proposition}
\begin{proof}
The number $\zeta$ is not a root of unity: $\zeta+\zeta^{-1}=6/5$
is rational but not an algebraic integer. The powers of $R_j$ and $D_j$
therefore have closures containing two distinct one-parameter subgroups
in the corresponding embedded $\SU(2)$. Their Lie brackets generate
that $\mathfrak{su}(2)$. The adjacent embedded Lie algebras generate
$\mathfrak{su}(q)$: adjacent elementary off-diagonal directions produce
all other off-diagonal directions by brackets, and their self-brackets
produce the traceless diagonal directions. The closed generated subgroup
is consequently $\SU(q)$.

The adjoint entries are algebraic. The spectral-gap theorem of
Benoist--de Saxc\'e \cite[Theorem 1.2]{BdS} applies to an adapted
algebraic measure on this compact connected group with simple Lie
algebra. Equivalently, one may use the algebraic unitary result of
Bourgain--Gamburd \cite{BG12}. The symmetric average has an upper gap;
weight $1/2$ at the identity makes its spectrum nonnegative. Its
mean-zero operator norm is therefore strictly below one. This is the
external full-group theorem, not a calculation in the defining
$q$-dimensional representation.
\end{proof}

\begin{proof}[Proof of Corollary~\ref{cor:projective39}]
Corollary~\ref{cor:conjugation} gives the lower order, for every fixed
positive calibration $\kappa$. Theorem~\ref{thm:projective-upper} gives
the matching exact upper order. Proposition~\ref{prop:algebraic-gates}
supplies fixed finite alphabets. The quantum assertion follows from the
next proposition. Taking binary logarithms gives the label-bit formula.
\end{proof}

\begin{proposition}[A fixed quantum realization]\label{prop:quantum}
The projective experiment has an exact quantum realization on $\C^q$
using fixed seed resets, the prescribed unitary commands and fixed
query measurements. It is independent of $N$.
\end{proposition}
\begin{proof}
On seed $P$ reset the density matrix to
\[
 \sigma_P=I/q+\rho Y(P)=(1-\rho/r_q)I/q+(\rho/r_q)P.
\]
This is positive and has trace one, since $0<\rho\le1/10<r_q$.
Commands act by conjugation. For query $j$, use effects
$(I+F_j)/2,(I-F_j)/2$. Hilbert--Schmidt normalization implies
$\norm{F_j}_{\mathrm{op}}\le1$, so these are positive effects summing
to $I$. Their outcome mean is
$\tr(F_j\sigma_{U_wPU_w^*})=\rho\tr(F_jY(U_wPU_w^*))$.
This is precisely \eqref{eq:law}. The seed reset is a channel, not a
unitary; only one query is measured.
\end{proof}
The comparison is a fixed $q$-dimensional quantum register versus
$(q-1)\log_2N+O(1)$ classical \emph{label bits}. Dimension or entropy
advantages for quantum models are prior phenomena \cite{Gu,Thompson,Ghafari}.
The assertion here is the proved classical width order for this complete
worst-word numerical behavior. It does not request jointly independent
answers to all queries or optimize arbitrary adaptive experiments.
